\documentclass[aspectratio=169,11pt]{beamer}
\usepackage[T1]{fontenc}
\usepackage{lmodern}
\usepackage{listings}
\usepackage{booktabs}
\definecolor{navy}{RGB}{25,48,76}
\definecolor{teal}{RGB}{0,112,120}
\definecolor{codebg}{RGB}{246,247,249}
\setbeamercolor{structure}{fg=teal}
\setbeamercolor{frametitle}{fg=navy}
\setbeamercolor{title}{fg=navy}
\setbeamercolor{normal text}{fg=navy,bg=white}
\setbeamertemplate{navigation symbols}{}
\setbeamertemplate{footline}{\hfill\usebeamercolor[fg]{structure}\insertframenumber\,/\,\inserttotalframenumber\hspace{5mm}\vspace{3mm}}
\setbeamertemplate{itemize items}[circle]
\setbeamersize{text margin left=9mm,text margin right=9mm}
\lstset{language=Python,basicstyle=\ttfamily\small,keywordstyle=\color{teal}\bfseries,commentstyle=\color{gray},stringstyle=\color{teal},backgroundcolor=\color{codebg},frame=none,showstringspaces=false,columns=fullflexible,keepspaces=true,breaklines=true,xleftmargin=2mm,xrightmargin=2mm,aboveskip=8pt,belowskip=8pt}
\newcommand{\code}[1]{\texttt{#1}}
\title{The Visitor Pattern}
\subtitle{A lambda-calculus interpreter in Python}
\author{CS413}
\date{Fall 2026}
\begin{document}
\begin{frame}\titlepage\end{frame}

\begin{frame}{Why use a visitor?}
An expression tree supports several operations:
\begin{itemize}
\item Evaluate the expression to obtain a value.
\item Compute its free variables.
\item Other possible operations include printing or type checking.
\end{itemize}
\medskip
Each operation needs a case for every kind of expression.
\medskip

\alert{Visitor groups the cases for one operation in one object.}
The expression classes store syntax and provide an \code{accept} method.
\end{frame}

\begin{frame}{The roles in the Visitor Pattern}
\renewcommand{\arraystretch}{1.5}
\begin{tabular}{p{.25\textwidth}p{.65\textwidth}}
\toprule
Role & In this interpreter \\
\midrule
Element & An expression node, such as \code{D0Eint} \\
\code{accept(visitor)} & Calls the visitor method for that node's form \\
Visitor interface & \code{D0ExpVisitor[R]} declares one method per form \\
Concrete visitor & \code{EvaluateVisitor} or \code{FreeVariableVisitor} \\
Client & Calls \code{expr.accept(visitor)} and uses the result \\
\bottomrule
\end{tabular}
\end{frame}

\begin{frame}[fragile]{A node delegates to the visitor}
\begin{lstlisting}
@dataclass
class D0Eint(D0E000):
    arg1: sint

    def accept[R](self, visitor: D0ExpVisitor[R]) -> R:
        return visitor.visit_int(self)
\end{lstlisting}
\begin{itemize}
\item The node knows that it is an integer expression.
\item It passes itself to \code{visit\_int}, so the visitor can read \code{arg1}.
\item The supplied visitor determines what to do with that integer.
\end{itemize}
\end{frame}

\begin{frame}[fragile]{One node, two operations}
\begin{lstlisting}
# In EvaluateVisitor:
def visit_int(self, dexp: D0Eint) -> d0val:
    return D0Vint(dexp.arg1)

# In FreeVariableVisitor:
def visit_int(self, dexp: D0Eint) -> frozenset[dvar]:
    return frozenset()
\end{lstlisting}
\begin{lstlisting}
expr = D0Eint(7)
expr.accept(EvaluateVisitor())      # D0Vint(7)
expr.accept(FreeVariableVisitor())  # frozenset()
\end{lstlisting}
The expression stays the same. The visitor selects the operation.
\end{frame}

\begin{frame}{How dispatch works}
For \code{expr = D0Eint(7)} and \code{v = EvaluateVisitor()}:
\begin{enumerate}
\item The client calls \code{expr.accept(v)}.
\item Python finds \code{accept} on \code{D0Eint}.
\item That method calls \code{v.visit\_int(expr)}.
\item Python finds \code{visit\_int} on \code{EvaluateVisitor}.
\item The visitor returns \code{D0Vint(7)} through \code{accept}.
\end{enumerate}
\medskip
The behavior depends on both the node class and the visitor class.
This is the Visitor Pattern's \alert{double dispatch}.
\medskip

Here, explicit method names such as \code{visit\_int} encode the node case.
\end{frame}

\begin{frame}[fragile]{The visitor controls recursive traversal}
\begin{lstlisting}
# In FreeVariableVisitor:
def visit_op2(self, dexp: D0Eop2) -> frozenset[dvar]:
    return dexp.arg1.accept(self) | dexp.arg2.accept(self)
\end{lstlisting}
For \code{D0Eop2("+", D0Evar("x"), D0Eint(1))}:
\begin{enumerate}
\item The root dispatches to \code{visit\_op2}.
\item The visitor visits the left child, obtaining \code{\{"x"\}}.
\item It visits the right child, obtaining the empty set.
\item It returns their union, \code{frozenset(\{"x"\})}.
\end{enumerate}
\medskip
\code{accept} dispatches one node. Visitor methods decide which children to visit.
\end{frame}

\begin{frame}{The extension tradeoff}
\textbf{Adding an operation}\par
Create a new visitor and implement its cases. Existing expression classes
and existing visitors can stay unchanged.
\bigskip

\textbf{Adding an expression form}\par
Add the node class and its \code{accept} method, extend the visitor interface,
and add the new case to every concrete visitor.
\bigskip

Visitor is useful when the node forms are relatively stable and new
operations are common.
\end{frame}

\begin{frame}[fragile]{The visitor result type}
\begin{lstlisting}
class D0ExpVisitor[R](ABC):
    @abstractmethod
    def visit_int(self, dexp: D0Eint) -> R:
        raise NotImplementedError
    # One abstract method for each expression form.
\end{lstlisting}
\begin{itemize}
\item \code{EvaluateVisitor}: \code{R = d0val}.
\item \code{FreeVariableVisitor}: \code{R = frozenset[dvar]}.
\item Each \code{accept[R]} has a local type parameter linking the visitor's result type to its own return type.
\end{itemize}
This source requires Python 3.12 or later. Annotations support static checking.
Abstract methods require concrete visitors to implement every case.
\end{frame}

\begin{frame}[fragile]{The interpreter's public operations}
\begin{lstlisting}
def d0exp_fvset(dexp: d0exp) -> frozenset[dvar]:
    return dexp.accept(FreeVariableVisitor())

def d0exp_evaluate(dexp: d0exp,
                   denv: d0env = ENVnil()) -> d0val:
    return dexp.accept(EvaluateVisitor(denv))
\end{lstlisting}
Expressions describe programs. Values describe results.
Environments bind names to values.
\medskip

The interpreter builds syntax with constructors. It does not include a parser.
\end{frame}

\begin{frame}[fragile]{Expression constructors}
\begin{itemize}
\item Constants and operators: \code{D0Eint}, \code{D0Ebtf}, \code{D0Eop1}, \code{D0Eop2}.
\item Variables and functions: \code{D0Evar}, \code{D0Elam}, \code{D0Efix}, \code{D0Eapp}.
\item Binding and control: \code{D0Elet}, \code{D0Eif0}.
\item Pairs: \code{D0Epair}, \code{D0Epfst}, \code{D0Epsnd}.
\end{itemize}
\begin{lstlisting}
D0Eop2("+", D0Eint(2), D0Eint(3))  # syntax for 2 + 3
D0Elam("x", D0Evar("x"))          # lambda x. x
D0Efix("f", "x", body)            # fix f(x). body
\end{lstlisting}
The meaning of \code{arg1}, \code{arg2}, and \code{arg3} depends on the constructor.
\end{frame}

\begin{frame}{Free variables and lexical binding}
A variable occurrence is free when no enclosing binder binds it.
\begin{align*}
\operatorname{FV}(x) &= \{x\} \\
\operatorname{FV}(42) &= \varnothing \\
\operatorname{FV}(\lambda x.\ e) &= \operatorname{FV}(e)\setminus\{x\} \\
\operatorname{FV}(\mathrm{fix}\ f(x).\ e) &= \operatorname{FV}(e)\setminus\{f,x\} \\
\operatorname{FV}(\mathrm{let}\ x=e_1\ \mathrm{in}\ e_2)
 &= \operatorname{FV}(e_1)\cup(\operatorname{FV}(e_2)\setminus\{x\})
\end{align*}
Operators, application, and pairs combine their children's free names.
A conditional includes names from its condition and both branches.
\end{frame}

\begin{frame}[fragile]{A let binder covers only its body}
\begin{center}\Large \code{let x = x + 1 in x + y}\end{center}
\begin{itemize}
\item Initializer: free names are $\{x\}$.
\item Body: free names are $\{x,y\}$ before accounting for the binder.
\item Remove the bound $x$ from the body's set, then take the union.
\end{itemize}
\[
\operatorname{FV} = \{x\}\cup(\{x,y\}\setminus\{x\}) = \{x,y\}
\]
\begin{lstlisting}
def visit_let(self, dexp):
    return (dexp.arg2.accept(self)
            | (dexp.arg3.accept(self) - {dexp.arg1}))
\end{lstlisting}
\end{frame}

\begin{frame}[fragile]{Environments and shadowing}
\code{ENVnil()} is empty. \code{ENVcns(name, value, tail)} adds a binding.
Lookup searches the newest binding first.
\begin{lstlisting}
old = ENVcns("x", D0Vint(1), ENVnil())
new = ENVcns("x", D0Vint(2), old)

d0env_search(new, "x")  # D0Vint(2)
d0env_search(old, "x")  # D0Vint(1)
\end{lstlisting}
Frozen environment nodes preserve earlier scopes when a new node extends them.
\medskip

An unbound name returns \code{D0V000()}, the implementation's error sentinel.
\end{frame}

\begin{frame}[fragile]{Evaluating let introduces a new scope}
\begin{lstlisting}
def visit_let(self, dexp: D0Elet) -> d0val:
    value = dexp.arg2.accept(self)
    denv = ENVcns(dexp.arg1, value, self.denv)
    return dexp.arg3.accept(EvaluateVisitor(denv))
\end{lstlisting}
\begin{enumerate}
\item Evaluate the initializer in the current environment.
\item Extend that environment with the new binding.
\item Evaluate the body using a new visitor for the extended scope.
\end{enumerate}
The initializer cannot use the new binding. The caller's scope stays intact.
\end{frame}

\begin{frame}[fragile]{Values and closures}
\code{D0Vint} and \code{D0Vbtf} hold scalar results.
\code{D0Vpair} holds two evaluated values.
\medskip

A \alert{closure} stores a function expression together with its defining environment.
\begin{lstlisting}
def visit_lam(self, dexp):
    return D0Vlam(self.denv, dexp)

def visit_fix(self, dexp):
    return D0Vfix(self.denv, dexp)
\end{lstlisting}
The body runs when the function is applied. This implementation captures the
whole environment, without trimming it using free-variable analysis.
\end{frame}

\begin{frame}[fragile]{Function application uses call by value}
\begin{lstlisting}
dfun = dexp.arg1.accept(self)  # function in caller's scope
darg = dexp.arg2.accept(self)  # argument in caller's scope

# When dfun is a D0Vlam:
dlam = dfun.arg2
denv = ENVcns(dlam.arg1, darg, dfun.arg1)
return dlam.arg2.accept(EvaluateVisitor(denv))
\end{lstlisting}
\begin{itemize}
\item Evaluate the function, then the argument, even if the body ignores it.
\item Bind the parameter in the \emph{closure's} environment.
\item Evaluate the body in that extended environment.
\end{itemize}
\end{frame}

\begin{frame}[fragile]{Lexical scoping preserves the defining environment}
Surface notation for a constructor-built expression:
\begin{lstlisting}
let x = 10 in
let f = (lam y. x + y) in
let x = 100 in
f(1)
\end{lstlisting}
\begin{enumerate}
\item Defining \code{f} captures the environment where \code{x = 10}.
\item Calling \code{f(1)} binds \code{y = 1} in that saved environment.
\item The body finds \code{x = 10} and \code{y = 1}.
\end{enumerate}
\medskip
Result: \alert{\code{D0Vint(11)}}. The caller's \code{x = 100} does not replace the captured binding.
\end{frame}

\begin{frame}[fragile]{Recursive closures bind their own name}
For \code{fix fact(n). if n == 0 then 1 else n * fact(n - 1)},
applying the resulting closure to $3$ produces \code{D0Vint(6)}.
\begin{lstlisting}
# When dfun is a D0Vfix:
dfix = dfun.arg2
denv = ENVcns(dfix.arg1, dfun, dfun.arg1)  # function name
denv = ENVcns(dfix.arg2, darg, denv)       # parameter
return dfix.arg3.accept(EvaluateVisitor(denv))
\end{lstlisting}
Each call extends the defining environment with the same recursive closure
and the current argument. Closure creation needs no cyclic environment.
\end{frame}

\begin{frame}[fragile]{Conditionals choose which branch to visit}
\begin{lstlisting}
def visit_if0(self, dexp: D0Eif0) -> d0val:
    cond = dexp.arg1.accept(self)
    if not isinstance(cond, D0Vbtf):
        raise TypeError(f"D0Vbtf(...) expected: {cond}")
    return (dexp.arg2 if cond.arg1 else dexp.arg3).accept(self)
\end{lstlisting}
Despite its name, \code{if0} requires a Boolean value in this interpreter.
\begin{lstlisting}
term = D0Eif0(D0Ebtf(True), D0Eint(7),
              D0Eop2("/", D0Eint(1), D0Eint(0)))
d0exp_evaluate(term)  # D0Vint(7)
\end{lstlisting}
The unselected branch never evaluates.
\end{frame}

\begin{frame}[fragile]{Pairs, operators, and runtime checks}
\begin{itemize}
\item Pairs evaluate both components, left to right. Projections require \code{D0Vpair}.
\item Unary operators \code{+1} and \code{-1} increment and decrement integers.
\item Binary arithmetic supports \code{+}, \code{-}, \code{*}, and \code{/}.
\item Comparisons \code{<}, \code{>}, \code{<=}, \code{>=}, \code{==}, and \code{!=} take integers and return \code{D0Vbtf}.
\end{itemize}
\begin{lstlisting}
d0exp_evaluate(D0Eop2("/", D0Eint(-3), D0Eint(2)))
# D0Vint(-2): division uses Python's floor division (//)
\end{lstlisting}
Wrong operand types raise \code{TypeError}. Division by zero raises
\code{ZeroDivisionError}. There is no static type checker here.
\end{frame}

\begin{frame}[fragile]{Analysis and evaluation visit different children}
\begin{lstlisting}
term = D0Eif0(D0Ebtf(True),
              D0Eint(7), D0Evar("missing"))

d0exp_fvset(term)     # frozenset({"missing"})
d0exp_evaluate(term)  # D0Vint(7)
\end{lstlisting}
\textbf{Free-variable analysis} visits all branches because it describes syntax.
\medskip

\textbf{Evaluation} visits the selected branch because it follows execution.
\medskip

Both use the same \code{accept} protocol. Each visitor defines its own traversal.
\end{frame}

\begin{frame}{Questions for discussion}
\begin{enumerate}
\item What is $\operatorname{FV}(\lambda x.\ x+y)$?
\medskip
\item What happens when we evaluate $(\lambda x.\ 42)\,(1/0)$?
\medskip
\item Why does the lexical-scope example return $11$?
\medskip
\item Which classes change if we add a \code{D0Eneg} expression?
\medskip
\item Does \code{accept} itself recursively visit every child?
\end{enumerate}
\end{frame}

\begin{frame}{Discussion answers}
\begin{enumerate}
\item $\{y\}$, because the lambda binds $x$.
\medskip
\item Argument evaluation raises \code{ZeroDivisionError} before the body runs.
\medskip
\item The closure captures $x=10$, and application adds $y=1$.
\medskip
\item Add \code{D0Eneg.accept}, declare \code{visit\_neg} in \code{D0ExpVisitor}, and implement it in each concrete visitor.
\medskip
\item No. Each visitor method decides how to visit children.
\end{enumerate}
\end{frame}

\begin{frame}{Reading the implementation}
Source: \code{../lambda1\_vp.py}
\begin{enumerate}
\item Expression dataclasses and their \code{accept} methods.
\item \code{D0ExpVisitor[R]} and its abstract cases.
\item \code{FreeVariableVisitor} and the \code{d0exp\_fvset} wrapper.
\item Runtime values, environments, and \code{d0env\_search}.
\item \code{EvaluateVisitor} and the \code{d0exp\_evaluate} wrapper.
\end{enumerate}
\medskip
The node supplies its expression form. The visitor supplies the operation.
\end{frame}
\end{document}
